\section{Related work and boundary}

The mapped compiler families sit inside well-established quantum-compilation traditions. The TARE case descends from a stabilizer-formalism Pauli block-encoding construction \cite{schillo2026blockencoding}, while anticommuting-unitary partitioning predates the present regime analysis \cite{izmaylov2020unitary}. Pauli-level compiler optimization and global binary-symplectic simplification are also established directions \cite{paykin2023pcoast,yang2026symphony}. The StabPrep transfer relies on the standard stabilizer-circuit setting in which efficient tableau simulation and canonical forms are classical results \cite{aaronson2004stabilizer}. Dynamic programming, Pauli/symplectic representations, stabilizer synthesis, and family-specific donor constructions therefore receive no novelty credit merely for appearing in this programme.

The bounded candidate residual is methodological and cross-family: making donor regions, exact trade refutations, sufficiency/tightness, structural recognizability, and prospective prediction into separately receipted components that can succeed or fail independently. The current manuscript does not claim that no earlier work maps any compiler family in this way. A fresh literature closure must compare the exact composition against normal-form, exact synthesis, phase/regime, and certified resource-estimation literature before submission.

The StabPrep result is particularly important for novelty discipline because it prevents the paper from selling an artifact of the first two families as a universal principle. Any field-level statement must include the negative transfer.

\paragraph{Mapping a problem space into regions is an established methodology.}
The idea that a problem space should be partitioned into regions, rather than summarised by a best score on a benchmark, is not new. Instance Space Analysis, built on Rice's algorithm-selection framework, selects the instance features most correlated with algorithm performance, projects them into a low-dimensional space, and identifies regions in which each algorithm performs well or poorly; a new instance is then routed by the region it falls into \cite{christiansen2025instancespace}. Christiansen and Smith-Miles apply it to the quadratic assignment problem and use it for a second purpose that bears directly on the argument here: the projection exposes which parts of the instance space the standard benchmark library fails to cover, and they find that the established library has poor coverage of a large region. Their response is to construct new instance classes to fill it.

That is the same objection this programme raises against evaluating compilation by best cost on a fixed benchmark set, and it has been made more rigorously there than here. The difference lies in what the regions are regions \emph{of}. An instance space partitions by \emph{relative algorithm performance}: which method wins, and by how much, measured empirically. The geometry mapped here partitions by \emph{exactness}: whether a donor family is provably optimal on a region, which structural trades are refuted rather than merely outperformed, and how large a search family must be to be sufficient. A region where one metaheuristic reliably beats another is a different object from a region where a simple family is already exact and nothing can beat it, and only the second supports a refutation. The methodologies are complementary --- a regime map says where to stop searching, an instance space says what to run when searching is still required.

